%% ma: L10-B02-0036
%% lop: 10
%% chuong: 1
%% bai: 2
%% dang: 10.2.5
%% loai: DS
%% muc: [NB, TH, NB, NB]
%% dap_an: "ĐSSS"
%% nguon: "Ảnh giáo viên gửi 10/10/2026 (ThS. Nguyễn Hoàng Việt, Chương 2 Tập hợp), Đề 2 (đề thứ hai, không ghi tiêu đề), Phần 2 Câu 13"
%% kien_thuc: "Bài 2 (tập hợp và các phép toán trên tập hợp)"
%% trang_thai: cho-duyet
%% ly_do: "Dạng toán mới đề xuất 10.2.5: phép toán trên các tập hợp hữu hạn (hợp, giao, hiệu)"
\begin{ex}
Cho hai tập hợp: $A=\{x\in\mathbb R\mid(x-1)(x-2)(x-3)=0\}$; $B=\{5;3;1\}$. Vậy:
\choiceTF
{\True Tập hợp $A$ có $3$ phần tử.}
{Tập hợp $A\cup B$ có $6$ phần tử.}
{Tập hợp $A\subset B$.}
{Tập hợp $B\subset A$.}
\loigiai{$A=\{1;2;3\}$. a) Đúng. b) $A\cup B=\{1;2;3;5\}$ có $4$ phần tử, sai. c) $2\in A$ nhưng $2\notin B$ nên sai. d) $5\in B$ nhưng $5\notin A$ nên sai.}
\end{ex}
